\documentclass[11pt,a4paper]{article}

\usepackage[utf8]{inputenc}
\usepackage[T1]{fontenc}
\usepackage[italian]{babel}
\usepackage{geometry}
\usepackage{array}
\usepackage{xcolor}
\usepackage{helvet}

\renewcommand{\familydefault}{\sfdefault}
\geometry{margin=2.2cm}
\pagestyle{empty}
\setlength{\parindent}{0pt}
\setlength{\arrayrulewidth}{0.5pt}
\renewcommand{\arraystretch}{1.38}

\definecolor{accent}{HTML}{1F4E79}
\definecolor{lightgray}{HTML}{F3F6FA}

\begin{document}

\begin{center}
    {\Large\bfseries Dichiarazione progetto}\\[0.25cm]
    {\normalsize Tecnologie Web}
\end{center}

\vspace{0.55cm}

\begin{tabular}{>{\bfseries}p{3.75cm} p{11.05cm}}
Nome & Mattia \\
Cognome & Lemma \\
Matricola & N86005170 \\
Traccia scelta & REGEXLAB - piattaforma web per enigmi testuali basati su espressioni regolari. L'applicazione permette di proporre sfide, validare risposte tramite regex e gestire l'interazione tra utente e API. \\
Back-end & Django come framework principale, Django REST Framework per esporre API REST, Simple JWT per l'autenticazione, PostgreSQL come database relazionale, Gunicorn per l'esecuzione in produzione e drf-spectacular per la documentazione delle API. \\
Front-end & React con TypeScript per l'interfaccia utente, Vite per sviluppo e build, Tailwind CSS per lo stile, React Router per la navigazione, TanStack Query e Axios per la comunicazione con le API. \\
Infrastruttura & Docker e Docker Compose per avviare l'intero ambiente in modo riproducibile; nginx come reverse proxy nello stack orientato alla produzione. \\
\end{tabular}

\vfill

\begin{center}
    \colorbox{lightgray}{%
        \parbox{0.9\textwidth}{%
            \small
            Il progetto e' stato sviluppato separando back-end e front-end: il back-end espone API REST
            con Django e gestisce dati, autenticazione e logica applicativa, mentre il front-end in React
            e TypeScript costruisce l'interfaccia utente e comunica con le API. Docker Compose permette
            di avviare in modo coerente tutti i servizi necessari, inclusi database e reverse proxy.
        }%
    }
\end{center}

\end{document}
